\chapter*{APPENDIX}
\addcontentsline{toc}{chapter}{APPENDIX}
\label{chap:appendix}

This appendix provides supplementary reference material to support the reproducibility of the experiments and the deployment of the framework:
\begin{enumerate}
    \item \textbf{Section~\ref{app:features}:} Complete inventory of TON\_IoT network features, detailing preprocessing rules, imputation strategies, and encoded dimensions.
    \item \textbf{Section~\ref{app:hyperparams}:} Complete specification of hyperparameters and environmental configurations for the machine learning models and buffer managers.
    \item \textbf{Section~\ref{app:metric_formulations}:} Formal mathematical definitions and operational justifications for all continual learning and edge viability evaluation metrics.
    \item \textbf{Section~\ref{app:code_structure}:} Architectural layout of the software codebase and instructions for end-to-end execution.
\end{enumerate}

\section{Telemetry Attribute Dictionary and Preprocessing Specification}
\label{app:features}

Table~\ref{tab:feature_dictionary} catalogues the network attributes from the raw TON\_IoT telemetry dataset (\texttt{train\_test\_network.csv}), detailing their data types, roles, handling under Module~M1 preprocessing, and representation after train-only transformation.

\begin{table}[htpb]
\centering
\small
\renewcommand{\arraystretch}{1.3}
\caption{TON\_IoT telemetry feature inventory and preprocessing transformations}
\label{tab:feature_dictionary}
\begin{tabular}{|p{3.2cm}|p{2.2cm}|p{2.8cm}|p{4.6cm}|}
\hline
\textbf{Raw Field Name} & \textbf{Data Type} & \textbf{Preprocessing Action} & \textbf{Encoded Representation} \\ \hline
\texttt{type} & Categorical & Target attribute & Label encoded $\{0, \dots, K-1\}$; withheld attack assigned $K$ \\ \hline
\texttt{label} & Binary & Dropped & Redundant with multi-class \texttt{type} \\ \hline
\texttt{src\_ip}, \texttt{dst\_ip} & String / IP & Dropped & Stripped to prevent network memorisation \\ \hline
\texttt{src\_port}, \texttt{dst\_port} & Numeric (int) & Median imputed & Continuous numeric feature \\ \hline
\texttt{proto} & Categorical & Mode imputed & One-hot encoded (\texttt{tcp}, \texttt{udp}, \texttt{icmp}) \\ \hline
\texttt{service} & Categorical & Mode imputed & One-hot encoded (\texttt{http}, \texttt{dns}, \texttt{ssl}, etc.) \\ \hline
\texttt{duration} & Continuous & Median imputed & Flow duration (seconds) \\ \hline
\texttt{src\_bytes}, \texttt{dst\_bytes} & Numeric & Median imputed & Transferred volume in bytes \\ \hline
\texttt{conn\_state} & Categorical & Mode imputed & One-hot encoded connection status flags \\ \hline
\texttt{missed\_bytes} & Numeric & Median imputed & Dropped or unacknowledged bytes \\ \hline
\texttt{src\_pkts}, \texttt{dst\_pkts} & Numeric (int) & Median imputed & Packet counts per transmission direction \\ \hline
\texttt{src\_ip\_bytes}, \texttt{dst\_ip\_bytes} & Numeric & Median imputed & Total IP layer payload volume \\ \hline
\texttt{dns\_query}, \texttt{http\_uri} & String / Text & Pruned ($>50\%$ missing) & High-cardinality sparse fields removed \\ \hline
\end{tabular}
\end{table}

\section{System Hyperparameters and Operational Constants}
\label{app:hyperparams}

All experiments were executed with fixed deterministic seeds to guarantee strict experimental repeatability. Table~\ref{tab:appendix_hyperparams} provides the exact configuration parameters used across Modules~M1 through~M7.

\begin{table}[htpb]
\centering
\small
\renewcommand{\arraystretch}{1.3}
\caption{System parameters and algorithmic settings}
\label{tab:appendix_hyperparams}
\begin{tabular}{|p{4.0cm}|p{4.2cm}|p{4.6cm}|}
\hline
\textbf{Functional Module} & \textbf{Hyperparameter} & \textbf{Configured Value} \\ \hline
Global Environment & Random Seed ($s$) & $42$ (enforced across splits and models) \\ \hline
Data Partitioning (M1) & Known Train / Test Split & $70\% \text{ Train} / 30\% \text{ Test}$ (stratified) \\ \hline
Unseen Hold-out (M1) & Novel Adaptation / Test Split & $20\% \text{ Adaptation} (M) / 80\% \text{ Eval}$ \\ \hline
Random Forest Classifier (M2) & Ensemble Estimators ($B$) & $100\text{ decision trees}$ \\ 
& Maximum Tree Depth ($d_{\max}$) & $20\text{ levels}$ \\
& Splitting Criterion & Gini Impurity \\
& Minimum Samples per Split & $2$ \\ \hline
Isolation Forest Screener (M3) & Number of Trees & $100\text{ isolation trees}$ \\
& Subsample Size ($\psi$) & $256\text{ flows per tree}$ \\
& Contamination Parameter & \texttt{'auto'} (unsupervised, threshold $0.0$) \\ \hline
Dynamic Replay Manager (M5) & Retention Floor ($k_{\min}$) & $100\text{ flows per known class}$ \\
& Hard Storage Ceiling ($B_{\max}$) & $5{,}000\text{ flows total}$ ($k_{\max} = 555$) \\
& Volume Matching Factor ($\alpha$) & $1.0$ \\ \hline
Smart Dynamic Replay (M5) & Margin Uncertainty Ratio & Top $40\%$ boundary, $60\%$ uniform \\
& Effective Allocation & $\alpha = 1.1$, Cap $B_{\max} = 4{,}000$ flows \\ \hline
\end{tabular}
\end{table}

\section{Mathematical Formulations of Continual Learning Metrics}
\label{app:metric_formulations}

To evaluate model updates on an edge gateway across accuracy, stability, memory efficiency, and latency, the following metrics were used throughout this thesis.

\subsection{Per-Class Forgetting and Aggregate Forgetting}
Let $F1_{c}^{\text{before}}$ be the macro-F1 score on known class $c \in \{0, \dots, K-1\}$ evaluated prior to adaptation, and let $F1_{c}^{\text{after}}$ denote the performance of the adapted model on the same untouched evaluation set. The per-class catastrophic forgetting $\Delta_c$ is:
\begin{equation}
\Delta_c = F1_{c}^{\text{before}} - F1_{c}^{\text{after}}.
\end{equation}
Average Catastrophic Forgetting ($\mathcal{F}_{\text{avg}}$) across all $K$ known classes is defined as:
\begin{equation}
\mathcal{F}_{\text{avg}} = \frac{1}{K} \sum_{c=0}^{K-1} \Delta_c.
\end{equation}
Maximum Catastrophic Forgetting ($\mathcal{F}_{\max}$) measures the vulnerability of the single worst-hit class:
\begin{equation}
\mathcal{F}_{\max} = \max_{c \in \{0, \dots, K-1\}} \Delta_c.
\end{equation}
Backward Transfer ($\text{BWT}$) measures whether adaptation improved or degraded past knowledge:
\begin{equation}
\text{BWT} = -\mathcal{F}_{\text{avg}}.
\end{equation}

\subsection{Memory Reduction Ratio (MRR)}
The Memory Reduction Ratio measures the percentage of historical sample storage saved relative to storing the entire historical training archive $N = |\mathcal{D}_{\text{train}}|$:
\begin{equation}
\text{MRR} = 100 \times \left( 1 - \frac{|\mathcal{B}|}{N} \right).
\end{equation}

\subsection{Plasticity--Stability Index (PSI)}
The Plasticity--Stability Index computes the harmonic mean between new-task plasticity ($F1_{\text{new}}$) and historical stability ($1 - \max(0, \mathcal{F}_{\text{avg}})$):
\begin{equation}
\text{PSI} = \frac{2 \cdot F1_{\text{new}} \cdot (1 - \max(0, \mathcal{F}_{\text{avg}}))}{F1_{\text{new}} + (1 - \max(0, \mathcal{F}_{\text{avg}}))}.
\end{equation}
$\text{PSI} \in [0, 1]$, where $1.0$ indicates perfect zero-day acquisition with zero catastrophic forgetting.

\subsection{Edge IoT Viability Index (EIVI)}
The composite Edge IoT Viability Index combines detection quality, memory footprint, and adaptation latency into a single operational score scaled from $0$ to $100$:
\begin{equation}
\text{EIVI} = 100 \times \left[ 0.35\,F1_{\text{new}} + 0.35\,(1 - \max(0, \mathcal{F}_{\text{avg}})) + 0.25\,\left(\frac{\text{MRR}}{100}\right) + 0.05\,\min\!\left(1, \frac{1}{10\,\max(0.01, \Delta t)}\right) \right],
\end{equation}
where $\Delta t$ is the wall-clock retraining time in seconds.

\section{Software Repository Layout and Reproducibility Instructions}
\label{app:code_structure}

The experimental framework was implemented in Python using scikit-learn, pandas, and numpy. The code structure is organised as follows:

\begin{verbatim}
imple/
|-- configs/
|   `-- config.yaml                 # Master configuration and hyperparameter file
|-- src/
|   |-- preprocessing.py            # M1: Deduplication, IP stripping, data cleaning
|   |-- models/
|   |   `-- random_forest.py        # M2: Random Forest training and inference
|   |-- continual/
|   |   |-- buffer_manager.py       # M5: Static, Dynamic, and Smart Replay buffers
|   |   `-- dynamic_replay.py       # M6: Continual retraining pipeline
|   `-- evaluation/
|       `-- continual_metrics.py    # M7: F1, forgetting, PSI, and EIVI calculators
|-- notebooks/
|   |-- 03_baseline.py              # Phase 1-4 baseline evaluation script
|   |-- 05_isolation_forest.py      # Phase 5 anomaly screening script
|   |-- 07_1_buffer_sensitivity.py  # Phase 7.1 memory sweep script
|   `-- 10_master_consolidation.py  # Phase 8-10 multi-attack evaluation driver
`-- results/
    |-- metrics/                    # JSON evaluation outputs
    `-- tables/                     # CSV summary tables
\end{verbatim}

\subsubsection{Execution procedure}
All phases can be reproduced sequentially from the root directory using the virtual environment:
\begin{enumerate}
    \item \textbf{Data preparation and baseline execution:}
    \begin{verbatim}
    python notebooks/03_baseline.py
    \end{verbatim}
    This executes Algorithm~4.1, removes duplicates, fits the clean baseline on nine known classes, evaluates against the withheld attack, and saves the $10 \times 10$ confusion matrix.
    \item \textbf{Unsupervised anomaly screening:}
    \begin{verbatim}
    python notebooks/05_isolation_forest.py
    \end{verbatim}
    This fits the Isolation Forest on known training traffic and computes flag rates across normal, known attack, and unseen partitions.
    \item \textbf{Replay buffer sensitivity analysis:}
    \begin{verbatim}
    python notebooks/07_1_buffer_sensitivity.py
    \end{verbatim}
    This sweeps per-class quotas $k \in \{25, 50, 100, 200, 500\}$ to map the Pareto knee.
    \item \textbf{Multi-attack Dynamic Replay evaluation:}
    \begin{verbatim}
    python notebooks/10_master_consolidation.py
    \end{verbatim}
    This runs all four leave-one-attack-out scenarios across all five update strategies, outputting the comprehensive results tables reported in Chapter~\ref{chap:results}.
\end{enumerate}

